\documentclass[a4paper,11pt]{article}
\usepackage[utf8]{inputenc}
\usepackage[italian]{babel}
\usepackage{amsmath,amssymb}
\usepackage{booktabs}
\usepackage{geometry}
\usepackage{hyperref}
\usepackage{listings}
\usepackage{xcolor}

\geometry{margin=2.5cm}

\lstset{
  basicstyle=\ttfamily\small,
  backgroundcolor=\color{gray!10},
  frame=single,
  breaklines=true
}

\title{Normalizzazione Rank-Based dei Timestamp\\per l'Esperimento del Problema~2}
\author{}
\date{}

\begin{document}
\maketitle

\section{Problema riscontrato}

L'esperimento del Problema~2 viene terminato dal kernel (OOM-kill) su tutti i dataset con timestamp in formato Unix epoch. Le cause identificate sono due, entrambe legate al fatto che il range temporale $T = t_{\max} - t_{\min}$ risulta enormemente pi\`u grande del numero di timestamp effettivamente distinti nel dataset.

\subsection{Causa 1: vettore della soluzione alternativa}

La classe \texttt{AlternativeSolution} alloca un vettore di $T - \mu + 1$ slot (uno per finestra temporale). Per dataset come \texttt{sx-stackoverflow}, con $T \approx 2.4 \times 10^8$ e $\mu \approx 1.2 \times 10^8$, questo vettore richiede circa 960\,MB solo per i puntatori \texttt{nullptr}.

\subsection{Causa 2: enumerazione esplicita di $W(v)$}

La funzione \texttt{compute\_W} enumera \emph{ogni singolo intero} nell'unione degli intervalli $[\lambda - \mu,\, \lambda]$. Se il range normalizzato \`e dell'ordine di milioni, $|W(v)|$ pu\`o raggiungere decine di milioni per un singolo nodo, e la \texttt{SketchCache} memorizza uno sketch per ciascuno di questi valori.

\section{Soluzione adottata: preprocessing di normalizzazione}

\textbf{Non \`e stato modificato alcun file dell'esperimento} (\texttt{esperimento.cpp}, \texttt{esperimento\_opt.cpp}, \texttt{plot\_analysis.py}, \texttt{process\_result.py}).

\`E stato creato un unico file aggiuntivo, \texttt{normalize\_timestamps.cpp}, che implementa una \textbf{normalizzazione rank-based} dei timestamp come step di preprocessing.

\subsection{Funzionamento}

Dato un dataset con archi nel formato \texttt{u v t}:
\begin{enumerate}
  \item \textbf{Passata~1}: raccoglie tutti i timestamp, li ordina, e rimuove i duplicati ottenendo la sequenza $t_0 < t_1 < \cdots < t_{D-1}$ di $D$ timestamp distinti.
  \item \textbf{Passata~2}: riscrive il file sostituendo ogni timestamp $t_i$ con il suo rango $i \in \{0, 1, \ldots, D-1\}$.
\end{enumerate}

Il file normalizzato viene salvato in \texttt{normalized\_datasets/}.

\subsection{Correttezza}

La normalizzazione rank-based preserva tutte le propriet\`a rilevanti per l'esperimento:
\begin{itemize}
  \item L'\textbf{ordinamento temporale} degli archi \`e invariante (il rango \`e una funzione monotona crescente del timestamp originale).
  \item Le \textbf{relazioni di vicinato} $N(v, t)$ restano identiche: l'insieme dei vicini di $v$ al timestamp $t$ non dipende dal valore numerico di $t$, ma solo da quali archi hanno quel timestamp.
  \item Gli \textbf{sketch MinHash} e i \textbf{RangeTree} producono esattamente gli stessi risultati, perch\'e dipendono unicamente dall'insieme dei vicini a ciascun timestamp, non dal valore assoluto del timestamp.
  \item La \textbf{struttura LSH} e il \textbf{pruning $W(v)$} operano sugli stessi insiemi di finestre temporali valide.
\end{itemize}

\subsection{Effetto sulla compressione}

\begin{table}[h]
\centering
\begin{tabular}{lrrrc}
\toprule
\textbf{Dataset} & \textbf{Range originale} & \textbf{Timestamp distinti} & \textbf{Compressione} & \textbf{$\mu$ consigliato} \\
\midrule
aves-sparrow     & 1            & 2         & $1\times$        & 1  \\
chess\_year      & 8            & 9         & $1\times$        & 4  \\
CollegeMsg       & 16\,736\,181  & 58\,911    & $284\times$      & 10 \\
email-Eu-core    & 69\,459\,254  & 207\,880   & $334\times$      & 5  \\
soc-flickr       & 17\,017\,200  & 134        & $127\,949\times$ & 10 \\
sx-stackoverflow & 239\,705\,551 & 52\,423\,237 & $4.6\times$    & -- \\
wiki-talk        & 200\,483\,882 & 7\,375\,042  & $27\times$     & -- \\
\bottomrule
\end{tabular}
\caption{Effetto della normalizzazione rank-based.}
\end{table}

\subsection{Calibrazione di $\mu$}

Dopo la normalizzazione, $\mu$ va espresso in termini di \emph{numero di timestamp distinti consecutivi} nella finestra. La memoria occupata scala linearmente con $\mu$: su una macchina con 8\,GB di RAM (di cui $\sim$4\,GB disponibili), il limite pratico \`e circa $\mu \leq 50$ per CollegeMsg.

I dataset \texttt{sx-stackoverflow} e \texttt{wiki-talk} hanno rispettivamente 52M e 7M timestamp distinti: anche dopo la normalizzazione, la \texttt{TemporalRangeForest} da sola (con 2.6M e 1.1M nodi) eccede la memoria disponibile. Per questi dataset servirebbero almeno 32\,GB di RAM.

\section{File creati}

\begin{table}[h]
\centering
\begin{tabular}{ll}
\toprule
\textbf{File} & \textbf{Descrizione} \\
\midrule
\texttt{normalize\_timestamps.cpp} & Normalizzatore rank-based dei timestamp \\
\texttt{normalized\_datasets/*}    & Dataset con timestamp normalizzati \\
\texttt{normalized\_datasets/normalization\_info.csv} & Metadati della normalizzazione \\
\texttt{run\_all.sh}               & Script bash per esecuzione completa \\
\bottomrule
\end{tabular}
\end{table}

\section{Riepilogo}

\begin{enumerate}
  \item \textbf{Nessun codice dell'esperimento \`e stato modificato.}
  \item \`E stato aggiunto un \textbf{preprocessing} che normalizza i timestamp da valori Unix epoch a ranghi consecutivi $\{0, 1, \ldots, D-1\}$.
  \item L'esperimento viene lanciato sui file normalizzati con valori di $\mu$ calibrati sul range compresso.
\end{enumerate}

\end{document}
